% ============================================================
%  Hans Brøchner, "On Søren Kierkegaard's Work as a Religious
%  Author"
%  Fædrelandet, vol. 16, Saturday 1 December 1855, no. 281
%  (signed "—r.").
%
%  TRANSLATION (English)
%  Translator: [HH]
%  Source: the 1855 article (public domain), as reprinted in
%  C.H. Koch (ed.), Brøchner og Kierkegaard (Scientia Danica,
%  Series H, Humanistica 8, vol. 13). Koch's editorial apparatus
%  is not translated; only Brøchner's article, his own footnote,
%  and his own parenthetical citations are rendered here.
%  Kierkegaard's work titles are given in their standard English
%  forms.
% ============================================================
\documentclass[12pt,a4paper]{article}
\usepackage[utf8]{inputenc}
\usepackage[T1]{fontenc}
\usepackage{libertinus}
\usepackage{libertinust1math}
\usepackage[english]{babel}
\usepackage[protrusion=true,expansion=true]{microtype}
\usepackage{setspace}
\usepackage[margin=1.2in]{geometry}
\usepackage{fancyhdr}
\usepackage{hyperref}

\hypersetup{
  pdftitle={On Søren Kierkegaard's Work as a Religious Author},
  pdfauthor={Hans Brøchner},
  hidelinks
}

\pagestyle{fancy}
\fancyhf{}
\rhead{\emph{Fædrelandet}, 1 Dec.\ 1855}
\cfoot{\thepage}

\setstretch{1.3}

\title{\textbf{On Søren Kierkegaard's Work\\ as a Religious Author}}
\author{[Hans Brøchner]}
\date{Translated from the Danish by [HH]\\[0.3em]
  \emph{Fædrelandet}, vol.\ 16, Saturday 1 December 1855, no.\ 281
  (signed ``---r.'')}

\begin{document}
\maketitle
\thispagestyle{fancy}
\bigskip

\noindent It could not have been otherwise than that the death of a man who had
awakened in so many serious thoughts, strong passions---those of love and of
hatred---must have set these thoughts, these passions, into living motion. At such
a moment it is natural to strive to survey what is now concluded, to seek in the
clear image of what he willed and accomplished a sure hold for one's thoughts and
for one's judgment. To give such an image---one that the author of these lines has
been unable to find in what has been written about Kierkegaard since his
death---this presentation is meant to contribute; it keeps throughout to his own
words, and will attempt to show, on the one hand, the inner connection and unity
in the whole series of his writings, and, on the other, how closely his last
appearance is bound up with his earlier activity, how definitely it was
foreshadowed therein, and with what necessary consistency it had, under the given
circumstances, to proceed from it.

The literary activity that is here to be discussed begins in February 1843 with
\emph{Either/Or}. Before this work Kierkegaard had indeed published two writings:
\emph{From the Papers of One Still Living, Published Against His Will by S.K.}
(1838), and the master's dissertation \emph{On the Concept of Irony, with
Continual Reference to Socrates} (1841). But although the first of these
writings---in its polemic against the merging of the single individual into the
impersonal mass, in its conception of the relation of the preceding generation to
the present, in its intimation of the religious goal (p.\ 30), and finally in much
of its form (especially the preface)---recalls his later works, it nevertheless
stands, by its accidental occasion and limited subject, isolated, outside the
coherent series; and although the second and incomparably more significant
writing, the treatise on irony, develops a concept of high importance for K.'s
conception of the stages of existence, and although (in its second part,
especially pp.\ 287--304) it gives the categorical determination of what in the
first part of \emph{Either/Or} appears in personal existence and in the second
part finds its judgment, it is nevertheless essentially a learned work, in which
the philological and historical apparatus occupies the preponderant place, and
what gives it a more general significance is again taken up in its place in the
later writings, notably in ``Concluding Unscientific Postscript.''

The properly coherent literary activity thus begins with \emph{Either/Or}, reaches
its decisive point in ``Concluding Unscientific Postscript'' (February 1846), its
provisional conclusion in ``Two Discourses at the Communion on Fridays'' (late
summer 1851; cf.\ the preface to these), and its final conclusion only at the
author's death. ``Concluding Postscript''---which, if any single one of K.'s
writings is to be singled out, must be designated his chief work---forms the
midpoint in the series of his writings and divides them into two groups, of which
the first strives toward the Christian, the second presents the Christian thus
attained.

The Christian is thus the goal, the task, and is so from the very first. ``The
authorship, regarded as a totality, is religious from first to last.'' ---
``The problem, the problem of the whole authorship, is: becoming a Christian.''
--- ``To this authorship there corresponds, as its source, one who \emph{qua}
author has willed only one thing.'' --- More precisely, this that he has willed is
determined thus: ``Without authority, to call attention to the religious, the
Christian.'' (On My Work as an Author, pp.\ 7, 9, 14). But those who were to be
made aware were those who in ``Christendom'' live in ``the sensory illusion of
calling themselves Christians, perhaps imagining themselves to be such, without
being so.'' The task was then not in the first instance to give, but first, as it
were, to take away, to remove what favored the illusion, and thus, by the
maieutic, by that ``art of the midwife,'' as Socrates called it, to help the
single individual himself ``to give birth to the truth.'' But the manner in which
the illusion was to be combated was conditioned by its ground. This lay partly in
the whole tendency of the age: that, reflective and dispassionate, it had let the
understanding level everything, and let ``the single individual'' disappear into
``the phantom of abstraction, the public''; that it, ``essentially sensible,
perhaps on the average more knowing than any generation before it,'' had
``forgotten what it means to exist and what inwardness is''; that it, impersonally
and disinterestedly losing itself in the objective, had forgotten appropriation,
forgotten the task of existing in one's knowledge; --- partly it lay in the
confusing intermixture of the different views of life, grounded in the fact that
the secularization of Christianity and the forgetting of the passion of
personality had led to the forgetting of Christianity's ideal requirement. To
dispel the illusion, these views of life had first to be brought forward in their
pronounced distinctiveness, and brought forward not as an object of
knowledge---for in that way the very tendency of the age away from existence and
from existence's decision would be favored---but in a form that asserts the
significance of existence, that compels the reader to acquire the understanding
personally, by himself. He is ``made aware''; the rest is left to himself; the
indirect communication conceals the communicator from him, makes him as it were
unrecognizable; he must now help himself in the appropriation, and in the
appropriation of that which, by being presented in existence, is only indirectly
communicated, he is personally to experience where the stumbling-block of the
lower view of life lies, where it strikes against passion, where existence loses
its foundation, and where despair leads either out into perdition or up to the
higher, to ``being saved infinitely in the essentiality of religiousness'' (cf.\ A
Literary Review. 1846. pp.\ 108 ff.). But since the earlier stages of life were to
be presented not by way of doctrine, as an object of knowledge, but in existing
individualities, the first series of writings had to be pseudonymous. The
pseudonymity had ``its essential ground in the production itself, which, for the
sake of the lines, for the sake of the psychologically varied differences of
individuality, poetically required that recklessness in good and evil, in
contrition and exuberance, in despair and arrogance, in suffering and jubilation,
and so on, that is ideally bounded only by psychological consistency, which no
factually real person, within reality's ethical limitation, dares to allow
himself, or could wish to allow himself'' (Appendix to the Concluding Postscript).
The pseudonyms are thus poetically created authors, like dramatic characters, who
live in their views and express these in their whole ideal consistency.

In the first of the pseudonymous writings, \emph{Either/Or, a Fragment of Life},
edited by Victor Eremita (1843), the two stages of life that lie prior to the
religious---the aesthetic and the ethical---are presented in existing
personalities. ``The aesthetic in a person is that by which he immediately is what
he is; whoever lives in and by and of and for this lives aesthetically''; his
watchword is enjoyment, his goal perdition. The first part presents such an
aesthetic individuality, tossed about among possibilities that his imagination
conjures up and his dialectic again annihilates, split apart in the passion of
isolated moments, without the coherence of existence, without actuality, captive
to despair. The second part presents a single individuality who, in despair's
relinquishment of the immediacy of the aesthetic, has made the movement of
infinity, whereby the ethical first comes into view; who in despair has found
himself again in his eternal validity, has chosen himself, who in this choice has
become manifest in the universal, and through this has gained the possibility of
giving his life continuity. The work's conclusion, an upbuilding discourse, leads
to the boundary of the religious, and, by bringing in the upbuilding, accentuates
the significance of subjectivity, of appropriation, in the relation to truth.

After \emph{Either/Or} there followed three smaller pseudonymous writings:
\emph{Fear and Trembling}, \emph{Repetition} (1843) and \emph{The Concept of
Anxiety} (1844), which successively bring forward the religious determinations,
whereas \emph{Either/Or} had given the aesthetic and ethical ones. They present in
existence the collision between the religious and the ethical, the ordeal, the
struggle in which the God-relationship is victorious, and finally the decisive
determination of the religious: sin, which posits the individual as heterogeneous
with the ethical, so that he cannot in despair assert himself by himself, win
himself out of despair, but only by a higher, by the God-relationship, in
self-annihilation. Simultaneously with the last of these writings there appeared a
little book, full of sparkling wit and biting satire: ``Prefaces, Light Reading
for Certain Classes as the Occasion May Require,'' which in its last section,
jestingly but in earnest, points to religious equality, to the requirement that
the highest for a human being must be that which is accessible to all, in contrast
to the accidental difference of talent. (``This thought of humanity and of human
equality: Christianly, every human being (the single individual), unconditionally
every human being, is equally near to God.'' (Preface to Two Discourses at the
Communion)).

The religious had now been reached in its essential determination, in the
determination of sin. This was not developed by way of doctrine---``for sin, like
faith, the paradox, and the like, has no place in the System''---but brought
forward in its psychological existence in the individual as anxiety. Now it was
certain that the Christian was not confused with anything lower, and
``Philosophical Fragments'' (1844) now set it forth in the form of a hypothesis,
experimentally. The form here conditioned pseudonymity, the subject-matter naming;
hence K.'s name appears on this writing as editor. The problem in this writing is:
``Can there be given a historical point of departure for an eternal consciousness;
how can such a one be of more than historical interest; can one build an eternal
happiness upon a historical knowledge?'' The hypothesis in it is: ``if one is to
go further than the Socratic,'' that is, than the view of life borne by the
thought of immanence, in which time, finitude, becomes the vanishing element
eternally taken up into the infinite; in which the settledness of eternity makes
existence, however earnestly it is taken, into a jest; in which the teacher
becomes merely the occasion for the learner, because the latter eternally has the
truth; in which the very consciousness of guilt does not reach beyond death,
because the finite personality does not reach beyond it. ``If, then, one is to go
further,'' other decisive determinations must be posited, and now, through a
construction that parodies speculation, the orthodox-Christian determinations come
forth in their whole simplicity. The hypothesis now becomes: ``a new organ:
faith; a new presupposition: the consciousness of sin; a new decision: the moment;
a new teacher: God in time.''

Simultaneously with these pseudonymous writings, K. had, under his own name,
published six small collections of upbuilding discourses, with constant aim at the
category of religion: ``the single individual,''\footnote{``It is Christian
heroism --- to venture wholly to become oneself, a single human being, this
definite single human being, alone before God, alone in this enormous exertion and
this enormous responsibility'' (The Sickness unto Death, Preface). --- It was to
this ``single individual, whom with joy and gratitude I call my reader,'' that the
prefaces to all the upbuilding discourses addressed themselves. On the relation of
the single individual to ``the congregation,'' K.\ has expressed himself more
fully in Practice in Christianity, pp.\ 235 f.\ (new edition); cf.\ On My Work as
an Author, p.\ 12, note; A Literary Review, pp.\ 107 f.} which drew nearer and
nearer to the Christian-religious, without, however, willing to reach it, since
they concluded in the humorous.

At the conclusion of Philosophical Fragments another writing had been intimated,
in which ``the problem was to be clothed in historical costume.'' But before this
there appeared ``Stages on Life's Way'' (1845), in which the three stages of
existence---the aesthetic, the ethical, and the religious---were more definitely
delineated, and the fundamental determination of the religious life, suffering,
was sharply accentuated. The year after there appeared ``Concluding Unscientific
Postscript to the Philosophical Fragments,'' like the ``Fragments'' with K.'s name
as editor. Here the problem of the subject's relation to Christianity, or:
becoming a Christian, is decisively posed; subjectivity as the truth for the
existing person is asserted against the objectivity-theories; the significance of
existence is asserted against speculation's looking-away from existence; indirectly
the point is indicated where the earlier views of life must run aground; the
spheres of existence are more definitely distinguished from one another; irony is
determined as the \emph{confinium} between the aesthetic and the ethical, humor as
the \emph{confinium} between the ethical and the religious; the religious sphere is
divided into two: the religiousness that rests in immanence, and the Christian,
whose fundamental determination---``in which all the other Christian determinations
lie, and from which they let themselves be consistently derived''---is this: ``that
in time the question of the individual's eternal happiness is decided, and decided
by the relation to Christianity as something historical,'' in which fundamental
determination the possibility of offense lies.

With this writing the goal, the Christian, had been decisively reached; it had
been brought forward in its absolute difference from the earlier stages, not as an
object of knowledge, but as the object of faith: the paradox, the despair of
knowledge, the stumbling-block of offense. The religious had been the goal; it had
``wholly and entirely been set into reflection, but in such a way that it was
wholly and entirely taken out of reflection back into simplicity''; the road
traversed had been ``to attain, to come to simplicity''; the task ``not to reflect
oneself into Christianity, but to reflect oneself out of something else and become,
more and more simply, a Christian.''

After the series of pseudonymous writings there followed the series of upbuilding
ones: first ``Upbuilding Discourses in Various Spirits'' (1847), in which, as the
title indicates, the upbuilding from the earlier stages was taken up, so that only
the last third gave the distinctively Christian; thereafter \emph{Works of Love,
Christian Deliberations in the Form of Discourses} (1847) and \emph{Christian
Discourses} (1848), to which in the immediately following years several smaller
collections attached themselves, ever more definitely striving toward what is the
midpoint of Christianity, so that, as K. then expressed it, the whole literary
activity might find ``its decisive resting-point at the foot of the altar, where
the author, himself best conscious of his own imperfection and guilt, by no means
calls himself a witness to the truth, but only a peculiar kind of poet and thinker
who, `without authority,' has had nothing new to bring, but has `wished to read
through the original text of the individual, humane conditions of existence---the
old, familiar text handed down from the fathers---once again, if possible in a
more inward way.''' (Preface to Two Discourses at the Communion on Fridays, 1851.)

The task of the upbuilding writings was ``with God's help to use everything in
order to make clear what Christianity's requirement in truth is,'' not to conduct
an apology for Christianity, but ``to present it, what truth is, as something so
infinitely high that the apology falls in another place, becomes, for us, that we
dare to call ourselves Christians, or transforms itself into a penitent
confession, that we thank God when we merely dare to regard ourselves as
Christians.'' But alongside the severity of the infinite requirement, the infinity
of grace was not forgotten; only the requirement had to ``be heard and asserted
first, heard and asserted in its whole infinity,'' before ``grace'' could be
applied. ``What I have willed is to contribute, by means of admissions, to
bringing if possible somewhat more truth into these (in the direction of being an
ethical and ethical-religious character, of renouncing worldly shrewdness, of
being willing to suffer for the truth, and the like) imperfect existences such as
we lead them, which is nevertheless always something, and in any case the first
condition for coming to exist more capably. What I have willed to prevent is that
one should then not, restricting oneself to and contenting oneself with the easier
and lower, thereupon go further, abolish the higher, --- go further, set the lower
in the higher's place, --- go further, make the higher into the fantast and
ridiculous exaggeration, the lower into wisdom and true earnestness; that one
should not then, in `Christendom,' existentially take Luther and the significance
of Luther's life in vain: this I have willed to contribute, if possible, to
preventing'' (On My Work as an Author, p.\ 18).

Christianity's requirement in its ``severity'' was asserted in two writings that
appeared pseudonymously, with the author's name as editor. In the first, \emph{The
Sickness unto Death, a Christian Psychological Exposition for Upbuilding and
Awakening} (1849), the sickness unto death---that is, despair---is conceived as
sin, and this as the opposite of faith, whose object is the paradox, the absurd
for thought; the possibility of offense is held fast as the dialectical moment in
all that is Christian. In the second writing, the one that marks the decisive
turning-point in K.'s relation to the existing official Church, \emph{Practice in
Christianity} (published 1850, written 1848), ``the requirement of being a
Christian was forced up to a highest of ideality'' in the direction of a relation
of contemporaneity to the prototype, of making the prototype's life present in our
own, of suffering, of self-annihilation, of faith in conflict with the
understanding under the inseparable possibility of offense. This book was a
question posed to the ecclesiastically established order. In the conception of
suffering as the inseparable mark of religiousness, in the conception of
Christianity as heterogeneity with the world, as a requirement of imitation of the
prototype, of ``transformation of existence in likeness to God,'' there was
pronounced a judgment upon a ``Christendom'' that has made itself comfortable in
the world, found rest in worldliness; upon a clergy that, protected by the State,
in the enjoyment of well-being proclaims the religion of suffering. What was later
expressed in ``The Moment'' was here not merely intimated but, in all essentials,
expressly stated (cf.\ especially pp.\ 41, 75, 123, 136, 209, 228, 241, 243 f.,
268 in the new edition). But, disinclined as K. was to work toward a change in
externals, he made, in the preface he added to the book, yet one more attempt to
give the established order a means of defense. ``My thought was: if this
established order is to be defensible, this is the only way: by poetically
(therefore by a pseudonym) lodging the judgment upon it, by then drawing in a
second power upon `grace,' so that it became Christianity, not merely by `grace' to
find forgiveness for the past, but by grace a kind of indulgence from the actual
imitation of Christ and the actual exertion involved in being a Christian. In this
way truth nevertheless enters into the established order; it defends itself by
judging itself, it concedes the Christian requirement, makes confession concerning
itself as to its distance, and without being able to be called a striving in the
direction of coming nearer to the requirement, but takes refuge in grace `also
with respect to the use one makes of grace.''' (``Fædrelandet'' for 16 May 1855;
cf.\ the Editor's Preface to Practice in Christianity, and: On My Work as an
Author, pp.\ 17, 19, note.)

The question was now posed to the official Church, in the first instance to its
head, the late Bishop Mynster; in the judgment upon it the head of the Church
would pronounce his own and the Church's own judgment. ``If there were strength in
him, he would have to do one of two things: either decidedly declare himself for
the book, dare to go along with it, let it be the defense that wards off what the
book poetically contains, the charge against the whole official Christianity that
it is an optical illusion, `not worth a rotten herring,' or else as decidedly as
possible throw himself against it, brand it as a blasphemous and sacrilegious
attempt, and declare the official Christianity to be the true Christianity.''
Neither was done: ``impotent, he chose to do nothing, except at most to condemn it
in the drawing-room; but not even in private conversation with me, although I asked
him to, after it had been brought before him, with his consent, what he was
judging in the drawing-room.'' (``Fædrel.'' in the place cited.)

So the matter stood while Mynster lived. Only the shyness about hastening a
decision in externals---a shyness grounded as much in K.'s conception of the
essence of the ethical and the religious as in his conception of the development
through which the age had to be healed (cf.\ the remarkable conclusion of ``A
Literary Review.'' 1846. pp.\ 107 ff.)---his piety toward Mynster, preserved from
the impressions of childhood, and probably also the hope, against all probability,
of seeing the thought ``so infinitely dear to him'' become reality, that Mynster
might still say: ``this is in truth Christianity, thus do I myself understand
Christianity in quiet inwardness,'' so that the book would become a transfiguration
of Bishop Mynster's proclamation of Christianity---could still hold him back from
the decisive step. Then Mynster died, and at his grave his successor proclaimed him
a witness to the truth, one of the genuine witnesses to the truth. With this
judgment the answer was given to the question that had been posed by ``Practice in
Christianity.'' The official Church had now declared that its proclamation of
Christianity, ``which toned down and kept silent the decisive Christian,'' was the
true one; that Christianity in its ideal requirement, the Christianity of the New
Testament, was exaggeration, overwroughtness. Now the protest had to be made
against this Church; now first the necessity of the opposition drove him to carry
the matter to the utmost (cf.\ ``This Must Be Said.'' p.\ 10). Yet it did not
happen rashly; only almost a year after Mynster's death and his successor's
canonizing of him was the protest made public (in ``Fædrel.'' for 18 December
1854). What has happened since that time is still in fresh memory. In the articles
in ``Fædrelandet'' and in ``The Moment'' the fight was waged for Christianity
against Christendom and its silent representative, the clergy. It was now declared:
``Whoever you are, whatever your life may otherwise be, my friend; --- by (if
otherwise you take part in it) refraining from taking part in the public worship of
God as it now is (with the claim of being the Christianity of the New Testament),
you have always one guilt the less, and a great one; you do not take part in making
a fool of God by calling that the Christianity of the New Testament which is not
the Christianity of the New Testament.'' (``This Must Be Said; So Let It Be Said.''
p.\ 3). The demand was made upon the State ``to take away that sensory illusion''
of a Church that in worldliness had been fused together with the worldly, of a
clergy ``whose existence, Christianly, is an untruth,'' ``whose Christian
earnestness finds its expression in this, that `witnesses to the truth'
make---what satirical self-contradiction!---career and fortune in the world by
depicting on Sundays how the truth must suffer in this world!''; whose whole
position compels it to conceal from the congregation what Christianity is, and
thus, so far as it depends on it, to make it impossible. All the contradictions
that lie in the clergy's position, all the deception it has brought into
``Christendom,'' were laid bare with biting satire: ``The Almighty knows best how
the blows are to be struck so that it is felt, that laughter, served in fear and
trembling, must be the scourge---therefore I am used.'' And against the wretched
half-measure that has retained merely the name of Christianity, against the
delusion ``that being a Christian is something everyone is as a matter of course,
something one comes by more easily than the most insignificant skill,'' the ideal
was set up, earnestly and in judgment. It is now two months since the last number
of ``The Moment'' appeared; a few days after, illness confined the author to the
sickbed from which he did not rise again. His activity is now concluded, and when
we look back over it, we shall find his words true, that he ``as an author has
willed only one thing.'' To call attention to the Christian was the task of his
life; it was the task in the first series of his writings, which worked against the
confusion of the Christian with what is not Christian, against the confounding of a
knowledge about Christianity with a transformation of existence by Christianity; it
was the task in his upbuilding writings, which, with deep knowledge of the hidden
passions of the human heart, knew how to track down every self-deception, to block
every way out by which one might evade becoming the Christian; it was the task in
his struggle against the Church that kept silent the Christian, ``made Christianity
into poetry, into something to which one relates oneself only through the
imagination,'' abolished the imitation of Christ in suffering, in dying away from
the world---the Church that ``cast the prototype away and let what one is oneself,
mediocrity, pass for the ideal.'' And what he, in his life so early ended, fought
for---with a strength of will that did not let itself be bent by the body's
weakness, and with a faith that did not despair---has not been without fruit. His
earlier writings have worked in quiet in an ever-larger circle; his ``Moments''
have, in their thousands, opened the eye and awakened serious thoughts that will
not be forgotten now that he is dead, but find new nourishment in the writings he
has left behind, in the life he has awakened in so many souls.

\bigskip
\noindent\hfill ---r.

\end{document}
